% Bogurodzica, stanzas 1-2: critical edition
% Build:  gregorio bogurodzica.gabc
%         lualatex bogurodzica-critical.tex   (twice; needs GregorioTeX 6 and Junicode)
\documentclass[11pt,a4paper]{article}
\usepackage[margin=2.2cm,top=2cm,bottom=2.4cm]{geometry}
\usepackage{fontspec}
\setmainfont{Junicode}[Path=/usr/share/fonts/opentype/junicode/,Extension=.otf,UprightFont=*-Regular,ItalicFont=*-Italic,BoldFont=*-SmBold,BoldItalicFont=*-SmBoldItalic,Numbers=OldStyle]
\usepackage{polyglossia}
\setdefaultlanguage[variant=british]{english}
\setotherlanguage{polish}
\usepackage{microtype}
\usepackage{gregoriotex}
\usepackage{booktabs,longtable,array}
\usepackage{enumitem}
\usepackage{multicol}
\usepackage[hidelinks]{hyperref}
\usepackage{xcolor}
\definecolor{rubric}{RGB}{160,30,30}

\setlength{\parindent}{0pt}
\setlength{\parskip}{0.55em}
\linespread{1.08}

\gresetinitiallines{1}
\grechangestyle{initial}{\fontsize{40}{40}\selectfont}
\grechangestaffsize{19}
\grechangedim{baselineskip}{60pt plus 5pt minus 5pt}{scalable}
\grechangestyle{annotation}{\small\color{rubric}}

\newcommand{\sig}[1]{\textup{\textsc{\MakeLowercase{#1}}}}
\newcommand{\pl}[1]{\foreignlanguage{polish}{\textit{#1}}}
\newcommand{\dipl}[1]{\foreignlanguage{polish}{\textup{#1}}}
\newcommand{\sect}[1]{\section*{#1}}

\title{\textsc{Bogurodzica}\\[0.3em]
  \large The two oldest stanzas with their melody\\[0.2em]
  \normalsize Critical edition}
\author{Edited from Kraków, Biblioteka Jagiellońska, MS 1619}
\date{2026}

\begin{document}
\maketitle
\thispagestyle{empty}

\begin{center}\small
\begin{minipage}{0.86\textwidth}
\emph{Bogurodzica} is the oldest surviving Polish religious song. Its first
two stanzas, addressed to the Virgin and to Christ, form the original song;
later copies add up to twenty more. The only medieval copy with music for both
stanzas is a flyleaf of BJ MS 1619, written shortly after 1400. This edition
gives that melody in square notation, with the text of the same copy, an
apparatus of musical and textual variants, and notes on the manuscript.
\end{minipage}
\end{center}

\sect{Edition}

\input{bogurodzica.gtex}

\vspace{0.4em}
{\footnotesize Bar lines are editorial. Neume groups follow \sig{K}.}

\newpage
\subsection*{Text}

\begin{multicols}{2}\small
\textbf{Diplomatic} (\sig{K}; \dipl{ǫ} = nasal vowel)\\[0.3em]
\dipl{Bogv rodzicza dzewicza bogem slawena maria\\
Vtwego syna gospodzina matko swolena maria\\
Sziszczi nam spwczi nam kyrieleyson}\\[0.4em]
\dipl{Twego dzela krzcziczela boszicze\\
Vslisz glosi naplen misli czlowecze\\
Slisz modlitwǫ yǫsz nosimi\\
A dacz raczi gegosz prosimi\\
a naswecze zbozni pobith\\
posziwocze raski pzbith ky'eleyson}

\columnbreak
\textbf{Critical}\\[0.3em]
\pl{Bogurodzica dziewica, Bogiem sławiena Maryja,\\
u twego syna Gospodzina matko zwolena, Maryja!\\
Zyszczy nam, spuści nam. Kirielejson.}\\[0.4em]
\pl{Twego dziela Krzciciela, Bożyce,\\
usłysz głosy, napełń myśli człowiecze.\\
Słysz modlitwę, jąż nosimy,\\
a dać raczy, jegoż prosimy:\\
a na świecie zbożny pobyt,\\
po żywocie rajski przebyt. Kirielejson.}
\end{multicols}

\textbf{Translation.} \textit{Mother of God, Virgin glorified by God, Mary!
Chosen Mother, Mary, win for us from thy Son the Lord, send down to us. Kyrie
eleison.}

\textit{For thy Baptist's sake, Son of God, hear our voices, fulfil the
thoughts of men. Hear the prayer we bring, and deign to grant what we ask: on
earth a godly life, after life a dwelling in paradise. Kyrie eleison.}

\sect{Sources}

\subsection*{Musical witnesses}
\begin{description}[leftmargin=2.6em,style=nextline,itemsep=0.3em,font=\normalfont]
\item[\sig{K}] Kraków, Biblioteka Jagiellońska, MS 1619, back flyleaf. Base
  text. Added to a copy of Latin sermons made by Maciej of Grochów, vicar of
  Kcynia; before 1408. Stanzas 1--2 with music, in seven systems. Gothic
  chant notation of German--Messine type, C and F clefs, no custodes. A
  decorated initial opens each musical section: \pl{B(ogu)}, \pl{U (twego)},
  \pl{Z(yszczy)}, \pl{T(wego)}, \pl{U(słysz)}, \pl{S(łysz)}, \pl{A (dać)}.
\item[\sig{K70}] Kraków, Biblioteka Jagiellońska, manuscript of 1570.
  Leans towards B\,flat before cadences on~a.
\item[\sig{W}] Warsaw, library of the Literati Chapel at St John's,
  MS 998\,D. Mensural notation, the only witness with written rhythm.
\item[\sig{Cz}] Częstochowa, Jasna Góra, late 15th century. Lost; known
  from a drawing reproduced by Woronczak. Music for the first two lines of
  stanza~1 only.
\item[\sig{Gn}] The Gniezno cathedral melody, printed in \emph{Przyjaciel
  Ludu} 4, no.\,46 (Leszno, 1838), p.~364. It sets the whole text, with
  smaller note values and sharps.
\item[\sig{Z}] Warsaw, Zamoyski Library, MS 1152 (1672). Mensural; the
  \sig{K70} version, less accurately written.
\end{description}
Readings of \sig{K70} and \sig{W} are given after Chybiński (1907,
pp.~30--41), who sets them out in modern notation. \sig{Cz}, \sig{Gn} and
\sig{Z} are described but not collated.

\subsection*{Text witnesses}
\begin{description}[leftmargin=2.6em,style=nextline,itemsep=0.2em,font=\normalfont]
\item[\sig{K408}] Kraków, BJ MS 408, f.~87\textsuperscript{v}. Added in the
  early 15th century to a codex dated 1408. Stanzas 1--14, no music.
\item[\sig{Wa}] Warsaw, Biblioteka Narodowa, from Holy Cross, c.\,1456.
  19 stanzas.
\item[\sig{Cz}] as above.
\item[\sig{K3}] Kraków, first half of the 16th century. 23 stanzas.
\item[\sig{Ł}] Jan Łaski, \emph{Commune incliti Poloniae Regni privilegium}
  (Kraków: J.\,Haller, 1506), first in the \emph{Registrum}, attributed to
  St Adalbert. No music.
\item[\sig{M}] Mateusz z Kościana, 1543. \quad \sig{H} Herburt, 1570.
  \quad \sig{S} Skarga, \emph{Żywoty świętych}, 1579.
\end{description}
Pilat (1879) and Kalina (1880) number the Kraków copies the other way round,
calling BJ 408 ``Kraków~I''. The sigla here follow the shelfmarks.

\sect{Editorial method}
\begin{enumerate}[itemsep=0.2em]
\item The melody and text are those of \sig{K}. The transcription was made
  from the facsimile and agrees in every pitch with Chybiński's (1907).
\item Pitch is fixed by the clefs of \sig{K}: final D, range C--d$'$,
  mode~1. With a C clef on the fourth line, gabc \texttt{c\,d\,e\,f\,g\,h\,i\,j\,k}
  equals C D E F G a b c$'$ d$'$. The single b, on \pl{spuś-ci}, is natural.
\item \sig{K} has no rhythmic signs, so the score has none.
\item The text is normalised after Woronczak (1962), keeping the \pl{u} that
  \sig{K} writes and sets to music at the start of line 1.2.
\item Bar lines mark the verse structure.
\end{enumerate}

\sect{Musical apparatus}

{\small Pitches C D E F G a b c$'$ d$'$; hyphens join notes on one syllable.
(?!) marks readings Chybiński thought corrupt.\par}

\renewcommand{\arraystretch}{1.25}
\begin{longtable}{@{}p{0.9cm}>{\raggedright\arraybackslash}p{3.2cm}>{\raggedright\arraybackslash}p{11.6cm}@{}}
\toprule
\textbf{Loc.} & \textbf{Lemma} & \textbf{Readings}\\
\midrule\endhead
1.1 & \pl{Bogurodzica} & \sig{K} D C F E-D-C D \quad \sig{K70} D D-C F E-C-E D \quad \sig{W} D C F E D\\
1.2 & \pl{dziewica} & \sig{K} F G-F-G a \quad \sig{K70} F G b\,flat a \quad \sig{W} F G a\\
1.3 & \pl{Bogiem sławiena} & \sig{K} c$'$ d$'$ c$'$ a-G-F F \quad \sig{K70} b d$'$ c$'$ a-G F (?!) \quad \sig{W} c$'$ d$'$ c$'$ a F\\
1.4 & \pl{Maryja} & \sig{K} F-G-a-G E-C D \quad \sig{K70} F-a-G-E C E D (?!) \quad \sig{W} F a G E D D (?!)\\
1.5 & \pl{U twego syna} & \sig{K} d$'$ c$'$ a-G-F G-F-G a \quad \sig{K70} a G F G b\,flat a a \quad \sig{W} short form beginning on a; \sig{K70} and \sig{W} both lack d$'$ c$'$\\
1.6 & \pl{Gospodzina} & \sig{K} c$'$ d$'$ c$'$ a G F \quad \sig{K70} c$'$ d$'$ c$'$ a-G F \quad \sig{W} c$'$ d$'$ c$'$ a F\\
1.7 & \pl{Matko} & \sig{K} c$'$ a-G-F \quad \sig{K70} c$'$ a-G (?!) \quad \sig{W} c$'$ a\\
1.8 & \pl{zwolena, Maryja} & as 1.4 in all witnesses\\
1.9 & \pl{Zyszczy nam} & \sig{K} F G-F-G a \quad \sig{K70} F G b\,flat a \quad \sig{W} F G a\\
1.10 & \pl{spuści nam} & \sig{K} \sig{W} c$'$ b d$'$ \quad \sig{K70} c$'$ b d$'$ d$'$ (?!)\\
1.11 & \pl{Kyrie} & \sig{K} c$'$ a G-F \quad \sig{K70} \sig{W} c$'$ a F\\
\midrule
2.1 & \pl{Twego dziela Krzciciela} & \sig{K} D-E D-C F G E C D \quad \sig{K70} D D-C F F G G E E C D D (?!) \quad \sig{W} D D-C F G E C-E D\\
2.2 & \pl{Bożyce} & as 1.4 in all witnesses\\
2.3 & \pl{usłysz głosy} & \sig{K} a G-F G-F-G a \quad \sig{K70} a G-F G b\,flat a \quad \sig{W} a G-F G a a\\
2.4 & \pl{napełń myśli człowiecze} & \sig{K} D C F G E C D \quad \sig{K70} \sig{W} \pl{napełnij}, with an extra note\\
2.5 & \pl{Słysz modlitwę, jąż nosimy} & \sig{K} C F G E-D-E C D E D \quad \sig{K70} C G a b\,flat a a \,\textbar\, D D C C F F E D \quad \sig{W} C F G a \,\textbar\, D D C F G a a\\
2.6 & \pl{a dać raczy} & \sig{K} a G-F G-F-G a \quad \sig{K70} a G F G b\,flat a a (?!)\\
2.7 & \pl{jegoż prosimy} & \sig{K} D-E D-C F G a \quad \sig{K70} D D-C F G b\,flat a\\
2.8 & \pl{a na świecie zbożny pobyt} & \sig{K} c$'$ a c$'$ G c$'$ a-G-F G-F-G a \quad \sig{K70} c$'$ b c$'$ d$'$ G G c$'$ a G G b\,flat a \quad \sig{W} a c$'$ d$'$ G G c$'$ a F G a a\\
2.9 & \pl{po żywocie rajski przebyt} & \sig{K} a c$'$ d$'$ G c$'$ a-G-F G-F-G a \quad \sig{K70} a b\,flat c$'$ d$'$ G \textbar\ c$'$ a F F b\,flat a \quad \sig{W} a c$'$ d$'$ G c$'$ a F G a a\\
2.10 & \pl{Kyrieleison} & \sig{K} a G F F-G-a-G E D D\\
\bottomrule
\end{longtable}

\sect{Textual apparatus}

{\small Variants are normalised; spelling differences are ignored.\par}

\begin{longtable}{@{}p{0.9cm}>{\raggedright\arraybackslash}p{3.3cm}>{\raggedright\arraybackslash}p{11.5cm}@{}}
\toprule
\textbf{Loc.} & \textbf{Lemma} & \textbf{Variants}\\
\midrule\endhead
1.1 & \pl{Bogu rodzica} & \sig{K} \sig{K408} \quad \pl{Boga rodzica} \sig{Wa} \sig{Cz} \sig{K3} \sig{Ł} \sig{M} \sig{H} \sig{S}\\
1.1 & \pl{sławiena} & \sig{K} \sig{K408} \quad \pl{sławiona} \sig{Wa} \sig{Cz} \sig{K3} \sig{Ł} \sig{M} \sig{H} \sig{S}\\
1.2 & \pl{U twego} & \sig{K} \sig{Wa} \sig{Cz} \sig{K3} \sig{Ł} \sig{H} \sig{S} \quad \pl{W twego} \sig{K408} \sig{M}\\
1.2 & \pl{Gospodzina} & \sig{K} \sig{K408} \sig{Wa} \sig{Cz} \sig{Ł} \sig{M} \quad \pl{gospodina} \sig{K3} \quad \pl{Gospodynia} \sig{H} \quad \pl{hospodyna} \sig{S}\\
1.2 & \pl{zwolena} & \sig{K} \sig{K408} \quad \pl{zwolona} \sig{Wa} \sig{Ł} \sig{M} \sig{H} \sig{S} \quad \pl{zbolenia} \sig{Cz} \quad \pl{zwolienia} \sig{K3}\\
1.3 & \pl{Zyszczy \dots\ spuści nam} & \sig{K} \sig{K408} \sig{Wa} \sig{Ł} \quad \pl{zyszczysz \dots\ spuszczysz} \sig{Cz} \quad \pl{zyszczy \dots\ spust winam} \sig{K3} \sig{M} \quad \pl{ziścisz \dots\ spuścisz} \sig{H} \quad \pl{ziści \dots\ spuści} \sig{S}\\
\midrule
2.1 & \pl{Twego dziela} & \sig{K} \sig{K408} \sig{Wa} \sig{Cz} \quad \pl{Twego syna} \sig{K3} \sig{Ł} \sig{M} \sig{H} \sig{S}\\
2.1 & \pl{Krzciciela} & \pl{zbawiciela} \sig{M} \quad \pl{[Kr]ziciela zbawiciela} \sig{K3}\\
2.1 & \pl{Bożyce} & \sig{K} \sig{K408} \quad \pl{zbożnica} \sig{Wa} \sig{Cz} \quad \pl{zbożny czas} \sig{Ł} \sig{H} \sig{S} \quad \pl{zbożnika} \sig{M}\\
2.2 & \pl{napełń} & \sig{K} \quad \pl{napełni} \sig{K408} \sig{Wa} \sig{Cz} \sig{K3} \sig{Ł} \sig{M} \sig{H} \sig{S}\\
2.3 & \pl{jąż nosimy} & \sig{K} \sig{K408} \sig{Wa} \quad \pl{jenżeć nosimy} \sig{Cz} \quad \pl{jenże (cię) prosimy} \sig{Ł} \sig{M} \sig{H} \sig{S}\\
2.4 & \pl{A dać raczy} & \sig{K} \sig{Wa} \quad \pl{Oddać raczy} \sig{K408} \sig{M} \quad \pl{O dać} \sig{Cz} \quad \pl{O o dać} \sig{Ł} \quad \pl{O o daci} \sig{H} \quad \pl{To dać} \sig{S}\\
2.5 & \pl{a na świecie} & \sig{K} \sig{K408} \quad \pl{daj na świecie} \sig{Wa} \sig{Cz} \sig{Ł} \sig{H} \sig{S} \quad \pl{daj nam na świecie} \sig{M} \quad \pl{dai na świecie dobry} \sig{K3}\\
2.5--6 & \pl{pobyt \dots\ przebyt} & \pl{pobytk \dots\ przybytk} \sig{Wa}\\
2.6 & \pl{rajski} & \dipl{raski} \sig{K}\\
refr. & \pl{Kirielejson} & \dipl{kyrieleyson} \sig{K} (st.\,2 \dipl{ky'eleyson}) \quad \dipl{kyon} \sig{Cz} \quad \dipl{kyeon} \sig{Wa} \quad \pl{Kyrie eleison} \sig{M} \sig{H} \sig{S}\\
\bottomrule
\end{longtable}

\sect{Notes}
\begin{description}[leftmargin=1.2em,itemsep=0.3em,font=\normalfont\itshape]
\item[Staves and clefs.] System 1 has four lines, the rest five; the lowest
  line also serves as the text baseline. In system~2 the clefs move for
  \pl{Matko} and return before \pl{zwolena}. Read with these clefs,
  \pl{Matko} is c$'$ a-G-F.
\item[1.1] \dipl{Bogv rodzicza} is two words, with the dative \pl{Bogu} in a
  genitive sense (Łoś). All other copies have \pl{Boga rodzica}, the name
  used in the 1410 \emph{Chronica conflictus} and by Długosz. \pl{Sławiena}
  and \pl{zwolena} are Old Polish forms, not Czechisms. The C on \pl{-gu-} is
  faint but legible.
\item[1.2] \pl{U}: two separate puncta, d$'$ c$'$, carry the syllable. The
  later musical witnesses drop both notes and so lose the octave leap from
  \pl{Maryja}. Over \pl{-dzi-na} of \pl{Gospodzina} the a and G stand close
  and the F apart, which gives c$'$ a-G \,\textbar\, F.
\item[2.1] \pl{dziela} means `for the sake of' (later \pl{dla}). Łoś knew it
  in no other Old Polish text and used it to date the song to the 13th
  century. The printed tradition replaced it with \pl{syna}; modern
  popular texts misread it as \pl{dzieła}. \pl{Bożyce} is the vocative of
  \pl{bożyc}, `son of God'.
\item[2.2] The C on \pl{-pełń} is not visible in the facsimile; only the D
  over \dipl{naplen} can be read. It is supplied from the parallel 2.1 and
  from every other transcription.
\item[2.4] The initial of \pl{A dać} has the same design as the \pl{U} of
  \pl{Usłysz}. This likely explains Pilat's reading \dipl{O dacz}; Łoś,
  Kalina and Woronczak read \dipl{A}.
\item[2.4--10] From \pl{A dać raczy} the neumes are larger and more widely
  spaced, possibly in another hand (Chybiński). The forms and pitches are
  the same as before.
\end{description}

\sect{Commentary}

\textbf{Form.} Each stanza has three parts: two parallel lines, then a
contrasting third with the refrain (Łoś). The melody varies four short
phrases (Surzyński's \emph{a--d}), all traceable to the first line, which
turns almost symmetrically on c$'$--d$'$--c$'$ (Flotzinger; Perz calls the
procedure centonate). The cadence of \pl{Maryja} returns at \pl{zwolena},
the second \pl{Maryja} and \pl{Bożyce}. \pl{Dziewica} returns at \pl{zyszczy
nam}, \pl{usłysz głosy} and \pl{a dać raczy}. \pl{Twego dziela Krzciciela}
returns at \pl{napełń myśli człowiecze}, and the end of \pl{a na świecie
zbożny pobyt} at \pl{po żywocie rajski przebyt}.

\textbf{Mode.} The final is D, with inner cadences on a and F. Chybiński
weighed mode~6 because of the F cadences, but settled on mode~1. B appears
once, at the melodic peak on \pl{spuści}.

\textbf{Date and origin.} Attribution to St Adalbert (d.\,997) begins with
15th- and 16th-century copies, Łaski among them, and no modern scholar
accepts it. Łoś placed the text in the late 13th century on linguistic and
metrical grounds; Chybiński and Feicht put the melody at the turn of the 13th
and 14th centuries. Woronczak (1952) matched the opening phrase to the Holy
Saturday litany \emph{Kyrie} in the Graz gradual MS 807 and read the song as
a Kyrie trope. Flotzinger (2005) showed that the match covers only the first
phrase, placed the Graz manuscript at Passau c.\,1170, and rejected the trope
reading. He dates the song to the late 14th century, in south-eastern Poland
with Czech mediation. Feicht (1962) assembled many parallels, among them
Jehan de Braine's \emph{Par desous l'ombre d'un bois}, and judged the melody
the coherent work of one trained musician. Pikulik (2005) argues that the
octave leap after the first \pl{Maryja} rules out a date before the 12th or
13th century and calls the melody a 13th-century cento. Kalina (1880) heard
two melodies by different authors; Feicht's analysis superseded this.

\sect{Bibliography}
{\small
\begin{description}[leftmargin=1.5em,itemsep=0.15em,font=\normalfont]
\item Adamiec, M., ed. \emph{Bogurodzica}. Wirtualna Biblioteka Literatury Polskiej, Uniwersytet Gdański. \url{https://literat.ug.edu.pl/bgrdzc/}.
\item Chmiel, A. ``Uwagi archiwalno-paleograficzne nad pieśnią \emph{Bogurodzica} w rękopisie Biblioteki Jagiell. nr 1619.'' \emph{Rozprawy Wydziału Filologicznego AU} 40 (1904).
\item Chybiński, A. \emph{``Bogurodzica'' pod względem historyczno-muzycznym}. Kraków, 1907.
\item Feicht, H. \emph{Studia nad muzyką polskiego średniowiecza}. Opera musicologica 1. Kraków: PWM, 1975.
\item Flotzinger, R. ``Zur \emph{Bogurodzica}-Frage.'' \emph{Context of Musicology} 1 (1997): 15--19. Polish trans. \emph{Pamiętnik Literacki} 96/2 (2005): 7--10.
\item Gołos, J. Review of Woronczak et al., \emph{Bogurodzica}. \emph{The Polish Review} 7/4 (1962): 84--85.
\item Kalina, A. \emph{Rozbiór krytyczny pieśni ``Bogarodzica''}. Lwów, 1880.
\item Łoś, J., ed. \emph{Bogurodzica. Pierwszy polski hymn narodowy}. Lublin, 1922.
\item Perz, M. ``Polskie `posłowie' do uwag Rudolfa Flotzingera.'' \emph{Pamiętnik Literacki} 96/2 (2005): 11.
\item Pikulik, J. ``Co melodia \emph{Bogurodzicy} mówi nam o czasie powstania pieśni.'' \emph{Pamiętnik Literacki} 96/2 (2005): 13--15.
\item Pilat, R. \emph{Pieśń ``Bogarodzica''. I. Restytucyja}. Kraków, 1879.
\item Surzyński, J. \emph{Polskie pieśni Kościoła katolickiego od najdawniejszych czasów do końca XVI stulecia}. Poznań, 1891.
\item Woronczak, J. ``Tropy i sekwencje w literaturze polskiej do połowy XVI wieku.'' \emph{Pamiętnik Literacki} 43 (1952): 353--54.
\item Woronczak, J., ed., with E. Ostrowska and H. Feicht. \emph{Bogurodzica}. Biblioteka Pisarzy Polskich A~1. Wrocław: Ossolineum, 1962.
\end{description}}

\end{document}
